% -------------------------------
% Plantilla del TFG en la ESIIAB-UCLM. Versión 0.0 (Preliminar)
%
% Luis de la Ossa
% 
% Compilar con XeLaTeX 
% -------------------------------


% -------------------------------
% Segunda portada
% -------------------------------
\cleardoublepage
\setcounter{page}{1} \null

\begin{textblock*}{3cm}(3cm,2cm) 
\begin{flushleft}
\includegraphics[height=2cm]{figs/logouclm.png}
\end{flushleft}
\end{textblock*}


\begin{textblock*}{2.5cm}(15.95cm,2.45cm) 
\begin{flushright}
\includegraphics[height=1.4cm]{figs/euroinf.jpg}
\end{flushright}
\end{textblock*}

\begin{textblock*}{2cm}(14cm,2cm)  % Corrige la posición porque la imagen tiene margen
\begin{flushright}
\includegraphics[height=2cm]{figs/logoesii.png}
\end{flushright}
\end{textblock*}



\begin{textblock*}{\textwidth}(3cm,8cm) 
\begin{center}\doublespacing 
{\fontsize{18pt}{4pt}\selectfont \bft{TRABAJO FIN DE GRADO}}

\medskip
{\fontsize{18pt}{4pt}\selectfont Grado en Ingeniería Informática}

\medskip
{\fontsize{16pt}{4pt}\selectfont \espec}
\end{center}
\end{textblock*}


\begin{textblock*}{\textwidth}(3cm,14cm) 
\begin{center}\setstretch{2}
{\fontsize{22pt}{4pt}\selectfont \bft{\titulo}}
\end{center}
\end{textblock*}


\begin{textblock*}{\textwidth}(3cm,20.5cm) 
\begin{flushleft}\doublespacing
{\fontsize{14pt}{4pt}\selectfont \bft{Autor:} \autor}

{\fontsize{14pt}{4pt}\selectfont \bft{Tutor:} \director}

{\fontsize{14pt}{4pt}\selectfont \bft{Co-Tutor:} \codirector}
\end{flushleft}
\end{textblock*}


\begin{textblock*}{\linewidth}(3cm,25cm) 
\begin{flushright}
{\fontsize{14pt}{4pt}\selectfont \fecha}
\end{flushright}
\end{textblock*}


% -------------------------------
% Declaración de autoría
% -------------------------------
\cleardoublepage
\thispagestyle{plain}
\begin{center}
\Large{\bft{Declaración de autoría}}
\end{center}
\vskip1cm

Yo, \autor, con DNI [NÚMERO DE DNI], declaro que soy el único autor del trabajo fin de grado titulado "Estudio comparativo de sistemas de ficheros
basados en Linux", que el citado trabajo no infringe las leyes en vigor sobre propiedad intelectual, y que todo el material no original contenido en dicho trabajo está apropiadamente atribuido a sus legítimos autores.

\vspace*{2cm}
\begin{center}
Albacete, a \quad \ldots \quad de \quad \ldots \quad de 20 \ldots

\vskip3cm

Fdo.: \autor
\end{center}


% -------------------------------
% Resumen
% -------------------------------
\cleardoublepage
\thispagestyle{plain}
\begin{center}
\Large{\bft{Resumen}}
\end{center}
\vskip1cm


La selección de un sistema de ficherosadecuado es crucial para lograr una gestión eficiente y efectiva de los datos en el contexto actual de la informática. Es en este marco que surge la idea de este proyecto, cuyo objetivo principal es llevar a cabo un análisis exhaustivo y multidimensional de los sistemas de ficheros relevantes para entornos Linux, abordando tanto sus aspectos teóricos como prácticos.\\

En la fase inicial, realizaremos una investigación minuciosa de los sistemas de ficheros ampliamente empleados en entornos Linux, destacando aquellos que han alcanzado un alto grado de madurez y aceptación. Este análisis se centrará en comprender las características fundamentales de estos Sistemas de Ficheros, evaluando sus ventajas y limitaciones, así como identificando los casos de uso para los cuales están diseñados de manera específica.\\

La siguiente etapa del proyecto implica la selección de un subconjunto de sistemas de ficheros para llevar a cabo un estudio más detallado y profundo. Esta fase de investigación abarcará evaluaciones teóricas de cada sistema. En la práctica, definiremos y ejecutaremos escenarios específicos para realizar pruebas detalladas, que incluirán la ejecución de \textit{benchmarks} para comparar el rendimiento relativo de cada sistema de ficherosen condiciones similares.\\

Además, estas pruebas se llevarán a cabo considerando la interoperabilidad a la que estos sistemas se enfrentan en situaciones del mundo real. Su objetivo es analizar el comportamiento de los sistemas frente a cargas de trabajo muy similares a las que se encontrarían en entornos reales, como entornos de usuario, de alto volumen de datos y concurrencia, o de alto trabajo de lectura y/o escritura.\\

Las conclusiones obtenidas a partir de este estudio proporcionarán una perspectiva informada y detallada para la selección del sistema de ficherosmás adecuado en entornos informáticos específicos. Se tendrán en cuenta no solo las características particulares y requisitos de las aplicaciones, sino también aspectos relacionados con el rendimiento, la escalabilidad y la mantenibilidad.\\

En última instancia, este enfoque integral y meticuloso busca contribuir de manera significativa al conocimiento práctico y a la toma de decisiones en la gestión de sistemas de ficheros para entornos Linux contemporáneos, fortaleciendo así la infraestructura y la eficacia tanto en entornos académicos como empresariales.\\

% -------------------------------
% Agradecimientos
% -------------------------------
\cleardoublepage
\thispagestyle{plain}
\begin{center}
\Large{\bft{Agradecimientos}}
\end{center}
\vskip1cm

